% Part: proof-theory
% Chapter: sequent-calculus
% Section: proof-examples

\documentclass[../../../include/open-logic-section]{subfiles}

\begin{document}

\olfileid{pt}{seq}{ex}

\olsection{د ثبوتونو بېلګې}

د سېکوېنټونو د !!{proof}s د ورکولو دپاره عموماً دا اسانه وي چې له وروستي سېکوېنټ څخه پورته خوا ته لاړ شو: په هغه کښې !!a{formula} د فعالې !!{formula} په توګه وټاکو، د اړوندې قاعدې مقدمې د سېکوېنټ د پاسه وليکو، او بيا دغه کړنه تر هغې تکرار کړو چې بديهي اصولو ته ورسېږو۔ د بېلګې په توګه، فرض کړئ غواړو دا سېکوېنټ ثابت کړو:
\[(!C \land !D) \lif !E \Sequent !C \lif (!D \lif !E)\]
$!C \lif (!D \lif !E)$ ته وګورئ: دا د $!A \lif !B$ په بڼه دے او د سېکوېنټ په ښۍ خوا کښې راځي۔ که ټول سېکوېنټ د \Log{G1c} له $\RightR{\lif}$ قاعدې سره برابر کړو،
\[\Axiom$ !A, \Gamma \fCenter \Delta, !B$
\RightLabel{\RightR{\lif}}
\UnaryInf$ \Gamma \fCenter \Delta, !A \lif !B $
\DisplayProof\]
نو راښيي چې بايد $\Gamma = \{(!C \land !D) \lif !E\}$، $\Delta =
\emptyset$، $!A \ident !C$ او $!B \ident !D \lif !E$ واخلو۔ په دې توګه د ثبوت وروستے استنتاج داسې کېدے شي:
\[\Axiom$ !C, (!C \land !D) \lif !E \fCenter !D \lif !E$
\RightLabel{\RightR{\lif}}
\UnaryInf$(!C \land !D) \lif !E \fCenter !C \lif (!D \lif !E)$
\DisplayProof\]
که همدغه لاره بيا وکاروو او $!D \lif !E$ د فعال فورمول په توګه وټاکو، نو دا ترلاسه کوو:
\[\Axiom$ !D, !C, (!C \land !D) \lif !E \fCenter !E$
\RightLabel{\RightR{\lif}}
\UnaryInf$!C, (!C \land !D) \lif !E \fCenter !D \lif !E$
\RightLabel{\RightR{\lif}}
\UnaryInf$(!C \land !D) \lif !E \fCenter !C \lif (!D \lif !E)$
\DisplayProof\]
اوس بايد $(!C \land !D) \lif !E$ ته پام وکړو او د \Log{G1c} $\LeftR{\lif}$ قاعده وکاروو:
\[\Axiom$ \Gamma \fCenter \Delta, !A$
\Axiom$ !B, \Gamma \fCenter \Delta$
\RightLabel{\LeftR{\lif}}
\BinaryInf$ !A \lif !B, \Gamma \fCenter \Delta$
\DisplayProof\]
له دې څخه دې بڼې ته رسېږو:
\[\Axiom$!D, !C \fCenter !E, !C \land !D$
\Axiom$!D, !C, !E \fCenter !E$
\RightLabel{\LeftR{\lif}}
\BinaryInf$ !D, !C, (!C \land !D) \lif !E \fCenter !E$
\RightLabel{\RightR{\lif}}
\UnaryInf$!C, (!C \land !D) \lif !E \fCenter !D \lif !E$
\RightLabel{\RightR{\lif}}
\UnaryInf$(!C \land !D) \lif !E \fCenter !C \lif (!D \lif !E)$
\DisplayProof\]
په کيڼ پورته سېکوېنټ کښې !!{formula} $!C \land !D$ د اصلي فورمول په توګه او $\RightR{\land}$ قاعده کارول دا راکوي:
\[
\Axiom$!D, !C \fCenter !E, !C$
\Axiom$!D, !C \fCenter !E, !D$
\RightLabel{\RightR{\land}}
\BinaryInf$!D, !C \fCenter !E, !C \land !D$
\Axiom$!D, !C, !E \fCenter !E$
\RightLabel{\LeftR{\lif}}
\BinaryInf$ !D, !C, (!C \land !D) \lif !E \fCenter !E$
\RightLabel{\RightR{\lif}}
\UnaryInf$!C, (!C \land !D) \lif !E \fCenter !D \lif !E$
\RightLabel{\RightR{\lif}}
\UnaryInf$(!C \land !D) \lif !E \fCenter !C \lif (!D \lif !E)$
\DisplayProof
\]
تر دې ځايه هر څۀ سم دي۔ پام وکړئ چې ټولې کارولې لارې په \Log{G3c} کښې هم کار کوي، ځکه تر دې ځايه کارولې قاعدې په \Log{G1c} او~\Log{G3c} کښې يو شان دي۔ د !!a{proof} داسې ټوټې ته رسېدلي يو چې په ټولو پورته سېکوېنټونو کښې هماغه !!{formula} په کيڼه او ښۍ خوا دواړو کښې شته۔ خو دا لا پوره بديهي اصول نۀ دي۔

په \Log{G1c} کښې بايد د ټولو پورته سېکوېنټونو !!{proof}s د $!A \Sequent !A$ په بڼه له بديهي اصولو څخه ورکړو۔ دا د کمزوري کولو د قواعدو په وسيله په اسانۍ کېږي؛ د بېلګې په توګه، تر ټولو کيڼ پورته سېکوېنټ $!D, !C \Sequent
!E, !C$ په لاندې ډول !!{prove} کېدے شي:
\[
\Axiom$!C \fCenter !C$
\RightLabel{\LeftR{\Weakening}}
\UnaryInf$!D, !C \fCenter !C$
\RightLabel{\RightR{\Weakening}}
\UnaryInf$!D, !C \fCenter !E, !C$
\DisplayProof
\]
په ټوليزه توګه، په \Log{G1c} کښې !!a{proof} لرو چې داسې ښکاري:
\[
\Axiom$!C \fCenter !C$
\RightLabel{\LeftR{\Weakening}}
\UnaryInf$!D, !C \fCenter !C$
\RightLabel{\RightR{\Weakening}}
\UnaryInf$!D, !C \fCenter !E, !C$
\Axiom$!D \fCenter !D$
\RightLabel{\LeftR{\Weakening}}
\UnaryInf$!D, !C \fCenter !D$
\RightLabel{\RightR{\Weakening}}
\UnaryInf$!D, !C \fCenter !E, !D$
\RightLabel{\RightR{\land}}
\BinaryInf$!D, !C \fCenter !E, !C \land !D$
\Axiom$!E \fCenter !E$
\RightLabel{\LeftR{\Weakening}}
\UnaryInf$!C, !E \fCenter !E$
\RightLabel{\LeftR{\Weakening}}
\UnaryInf$!D, !C, !E \fCenter !E$
\RightLabel{\LeftR{\lif}}
\BinaryInf$ !D, !C, (!C \land !D) \lif !E \fCenter !E$
\RightLabel{\RightR{\lif}}
\UnaryInf$!C, (!C \land !D) \lif !E \fCenter !D \lif !E$
\RightLabel{\RightR{\lif}}
\UnaryInf$(!C \land !D) \lif !E \fCenter !C \lif (!D \lif !E)$
\DisplayProof
\]
د دې ثبوت جوړښت، په تېره بيا لوړوالے ئې، په دې پورې نۀ تړلے چې !!{formula}s $!C$، $!D$ او~$!E$ څۀ دي۔ په پورته !!{proof} کښې دغه نښې په هر ډول !!{formula}s بدلولے شو، او پايله به بيا هم په~\Log{G1c} کښې !!{proof} وي۔ په \Log{G1c} کښې دا ثبوت \emph{سکيماتي} دے۔

په \Log{G3c} کښې د پېژندنې بديهي اصول د $!A, \Gamma \Sequent
\Delta, !A$ په بڼه سېکوېنټونه دي چې $!A$ بايد اتومي وي۔ نو که څۀ هم $!D, !C \Sequent
!E, !C$ د \Log{G3c} بديهي اصل ښکاري ($\Gamma = \{!D\}$؛ $\Delta =
\{!E\}$)، خو د پېژندنې دا بڼه هغه وخت \emph{په رښتيا} بديهي اصل ورکوي چې $!C$ په رښتيا اتومي وي۔ زموږ د ثبوت ټوټه په \Log{G3c} کښې بشپړ ثبوت دے چې $!C$، $!D$ او $!E$ اتومي !!{formula}s وي۔
د دې درو فورمولونو اتومي کېدل کافي دي، لازم شرط نۀ دے؛ د پاتې سياق يو ګډ اتومي فورمول هم هر پورته سېکوېنټ بديهي کولے شي۔ د کيڼې تشې نښې بديهي اصول هم په نظام کښې شته۔

که څۀ هم په کره توګه مو لا نۀ دي ښودلي چې
\[\Log{G3c} \Proves (!C \land !D) \lif !E \fCenter !C \lif (!D \lif
!E),\]
خو په عمل کښې همدومره کول راته بس دي او همدغه کولے شو۔ دا ځکه چې د $!A, \Gamma \Sequent \Delta, !A$ په بڼه هر سېکوېنټ په \Log{G3c} کښې د ثبوت وړ دے، که پخپله $!A$ اتومي هم نۀ وي۔ دا د~$!A$ پر ژوروالي~$\depth{!A}$ د استقرا په وسيله ښودلے شو۔

\begin{defn}\ollabel{defn:depth}
د !!a{formula} \emph{ژوروالے}~$\depth{!A}$ په لاندې ډول تعريفېږي:
\begin{enumerate}
  \item که $!A$ اتومي وي، نو $\depth{!A} = 0$۔
  \item که $!A \ident \lnot !B$، $!A \ident \lforall[x][!B]$ يا $!A \ident \lexists[x][!B]$ وي، نو $\depth{!A} = \depth{!B} + 1$۔
  \item که $!A \ident (!B \land !C)$، $(!B \lor !C)$ يا $(!B \lif
  !C)$ وي، نو $\depth{!A} = \max(\depth{!B},\depth{!C}) + 1$۔
\end{enumerate}
\end{defn}

دا ثابتول ګران نۀ دي چې د هر ترم~$t$ دپاره $\depth{A(x)} = \depth{A(t)}$ دے۔ زموږ دپاره مهمه خبره دا ده چې په هره قاعده کښې د اصلي !!{formula} ژوروالے تل د فعالو !!{formula}s له ژوروالي څخه زيات وي۔ د بېلګې په توګه، لرو:
\begin{align*}
\depth{P(x)} = \depth{P(a)} & = 0,\\
\depth{\lforall[x][P(x)]} & = 1,\\
\depth{(\lforall[x][P(x)] \lif P(a))} & = \max(1,0) + 1 = 2.
\end{align*}
له دې څخه پر~$\depth{!A}$ د استقرا په وسيله د ځينو پايلو د ثابتولو دپاره کار اخيستلے شو، لکه لاندې پايله۔

\begin{prop}\ollabel{prop:G1c-axioms}
د هر~$!A$ دپاره، په $\Log{G1c}$ کښې د $!A \Sequent !A$ !!a{proof} شته چې ټول بديهي اصول ئې اتومي دي۔
\end{prop}

\begin{proof}
  پر~$\depth{!A}$ د استقرا په وسيله۔ که $\depth{!A} = 0$ وي، نو $!A$ اتومي دے او $!A \Sequent !A$ له وړاندې په~\Log{G1c} کښې د غوښتل شوې بڼې !!{proof} دے۔ که $\depth{!A} > 0$ وي، نو $!A$ د لږ ژوروالي له !!{formula}s څخه جوړ شوے دے، لکه $!A \ident \lnot !B$، $!A \ident
  (!B \land !C)$ يا $!A \ident \lforall[x][!B(x)]$۔ د استقرا فرضيه دا ده چې د ټولو هغو $!D$ دپاره چې $\depth{!D} < \depth{!A}$ وي، $\Log{G1c} \Proves !D \Sequent !D$ له اتومي بديهي اصولو څخه حاصلېږي۔

  که $!A \ident \lnot !B$ وي، نو $\depth{!B} < \depth{!A}$ دے او د استقرا له فرضيې څخه په \Log{G1c} کښې د $!B \Sequent !B$ !!a{proof}~$\pi_1$ له اتومي بديهي اصولو څخه شته۔ دا د $\lnot !B \Sequent \lnot !B$ !!a{proof} ته په لاندې ډول غځولے شو:
  \[
  \AxiomC{}
  \RightLabel{$\pi_1$}
  \Deduce$!B \fCenter !B$
  \RightLabel{\LeftR{\lnot}}
  \UnaryInf$\lnot !B, !B \fCenter $
  \RightLabel{\RightR{\lnot}}
  \UnaryInf$\lnot !B \fCenter \lnot !B$
  \DisplayProof
\]

  که $!A \ident (!B \land !C)$ وي، نو $\depth{!B} < \depth{!A}$ او $\depth{!C} < \depth{!A}$ دي۔ د استقرا له فرضيې څخه په \Log{G1c} کښې د $!B
  \Sequent !B$ او $!C \Sequent !C$ !!{proof}s~$\pi_1$ او~$\pi_2$ له اتومي بديهي اصولو څخه شته۔ دا د $!B \land !C \Sequent !B \land !C$ !!a{proof} ته په لاندې ډول غځولے شو:
  \[
  \AxiomC{}
  \RightLabel{$\pi_1$}
  \Deduce$!B, \fCenter !B$
  \RightLabel{\LeftR{\Weakening}}
  \UnaryInf$!B, !C, \fCenter !B$
  \RightLabel{\LeftR{\land}}
  \UnaryInf$!B \land !C, !C\fCenter !B$
  \RightLabel{\LeftR{\land}}
  \UnaryInf$!B \land !C, !B \land !C \fCenter !B$
  \AxiomC{}
  \RightLabel{$\pi_2$}
  \Deduce$!C, \fCenter !C$
  \RightLabel{\LeftR{\Weakening}}
  \UnaryInf$!B, !C, \fCenter !C$
  \RightLabel{\LeftR{\land}}
  \UnaryInf$!B \land !C, !C \fCenter !C$
  \RightLabel{\LeftR{\land}}
  \UnaryInf$!B \land !C, !B \land !C \fCenter !C$
  \RightLabel{\RightR{\land}}
  \BinaryInf$!B \land !C, !B \land !C \fCenter !B \land !C$
  \RightLabel{\LeftR{\Contraction}}
  \UnaryInf$!B \land !C \fCenter !B \land !C$
  \DisplayProof
  \]

  اوس فرض کړئ $!A \ident \lforall[x][B(x)]$۔ داسې !!a{constant}~$c$ واخلئ چې په $\lforall[x][B(x)]$ کښې نۀ راځي۔ بيا $\depth{B(c)} =
  \depth{!B(x)} < \depth{!A}$ دے او د استقرا له فرضيې څخه په \Log{G1c} کښې د $!B(c)
  \Sequent !B(c)$ !!a{proof}~$\pi_1$ له اتومي بديهي اصولو څخه شته۔ دا د $\lforall[x][B(x)] \Sequent \lforall[x][B(x)]$ !!a{proof} ته په لاندې ډول غځولے شو:
  \[
  \AxiomC{}
  \RightLabel{$\pi_1$}
  \Deduce$!B(c) \fCenter !B(c)$
  \RightLabel{\LeftR{\lforall}}
  \UnaryInf$\lforall[x][!B(x)] \fCenter !B(c)$
  \RightLabel{\RightR{\lforall}}
  \UnaryInf$\lforall[x][!B(x)] \fCenter \lforall[x][!B(x)]$
  \DisplayProof
\]

  پاتې صورتونه د تمرينونو په توګه پرېښودل کېږي۔
\end{proof}

\begin{prob}
  د \olref{prop:G1c-axioms} ثبوت بشپړ کړئ: هغه صورتونه سر ته ورسوئ چې $!A \ident (!B \lor !C)$، $!A \ident (!B \lif !C)$ او $!A \ident \lexists[x][B(x)]$ وي۔
\end{prob}

زموږ د موخې دپاره، په هغه صورت کښې چې $!A \ident \lnot !B$ وي، دا !!{proof} هم کارولے شو:
\[
\AxiomC{}
\RightLabel{$\pi_1$}
\Deduce$!B \fCenter !B$
\RightLabel{\RightR{\lnot}}
\UnaryInf$\fCenter !B, \lnot !B$
\RightLabel{\LeftR{\lnot}}
\UnaryInf$\lnot !B \fCenter \lnot !B$
\DisplayProof
\]
خو داسې سېکوېنټي نظامونه هم شته چې دغه لاره پکې کار نۀ کوي۔ شهودي سېکوېنټي حساب \Log{G1i} د \Log{G1c} په شان دے، او ټول سېکوېنټونه داسې محدودوي چې په ښۍ خوا کښې ئې تر يوه زياته !!{formula} نۀ وي۔ که $\Delta = \emptyset$ وي، نو لومړنے !!{proof} به په~\Log{G1i} کښې هم !!a{proof} وي، خو دويم به نۀ وي، ځکه په يوه سېکوېنټ کښې $!B$ او~$\lnot !B$ دواړه په ښۍ خوا کښې دي۔
له دې سره د شرط د کيڼې قاعدې لومړۍ مقدمه هم بدلېږي: د هغې په ښي لوري کښې يوازې د شرط مقدم وي؛ د پاتې ښي سياق کاپي پکې نۀ راځي۔ يوازې د سېکوېنټونو د شمېر شرط د شهودي نظام د تعريف دپاره کافي نۀ دے۔

پام وکړئ چې ان په \Log{G1c} کښې هم، په هغه صورت کښې چې $!A \ident \lforall[x][!B(x)]$ وي، په عمومي ډول د قواعدو بل ترتيب نۀ شو کارولے۔ که ټاکلے ثابت په کيڼ فورمول کښې په رښتيا راځي، لاندې بڼه !!a{proof} \emph{نۀ ده}:
  \[
  \AxiomC{}
  \RightLabel{$\pi_1$}
  \Deduce$!B(c) \fCenter !B(c)$
  \RightLabel{\RightR{\lforall}}
  \UnaryInf$!B(c) \fCenter \lforall[x][!B(x)]$
  \RightLabel{\LeftR{\lforall}}
  \UnaryInf$\lforall[x][!B(x)] \fCenter \lforall[x][!B(x)]$
  \DisplayProof
\]
ځکه په~$\RightR{\lforall}$ کښې د ځانګړي متغير شرط ماتېږي۔
که کمي متغير په خپل بدن کښې هېڅ پېښه ونۀ لري، نو ټاکلے ثابت په کيڼ فورمول کښې هم نۀ راځي او د همدې ځانګړي متغير شرط نۀ ماتېږي؛ دا د پورته عمومي ممانعت استثنا ده۔

\begin{prop}\ollabel{prop:G3c-axioms}
د $!A, \Gamma
\Sequent \Delta, !A$ په بڼه هر سېکوېنټ (چې $!A$ ئې \emph{ضروري نۀ ده} اتومي وي) په~\Log{G3c} کښې د ثبوت وړ دے۔
\end{prop}

\begin{proof}
دا ثبوت د \olref{prop:G1c-axioms} ثبوت ته ډېر ورته دے، خو د استقرا په فرضيې کښې بايد احتياط وکړو۔ هغه صورت چې $n =
\depth{!A} = 0$ وي، بيا هم اسانه دے۔ که $n=0$ وي، نو $!A$ اتومي دے او $!A,
\Gamma \Sequent \Delta, !A$ د~\Log{G3c} بديهي اصل دے۔

اوس فرض کړئ $n>0$۔ بيا بېل بېل صورتونه څېړو۔ د استقرا فرضيه دا ده چې د ټولو هغو $!D$ دپاره چې $\depth{!D} < n$ وي، او د ټولو $\Pi$ او $\Lambda$ دپاره، $\Log{G3c} \Proves !D, \Pi \Sequent \Lambda, !D$ دے۔ دلته د $!A \ident (!B \land !C)$ صورت ښيو۔ د استقرا فرضيه دوه ځله کاروو: يو ځل د $!D = !B$، $\Pi = !C, \Gamma$ او $\Lambda
= \Delta$ دپاره، څو د $!B, !C, \Gamma \Sequent
\Delta, !B$ !!a{proof}~$\pi_1$ ترلاسه کړو؛ او بل ځل د $!D = !C$، $\Pi = !B, \Gamma$ او $\Lambda
= \Delta$ دپاره، څو د $!B, !C, \Gamma \Sequent
\Delta, !C$ !!a{proof}~$\pi_2$ ترلاسه کړو۔ دلته د~$!B \land !C, \Gamma \Sequent
\Delta, !B \land !C$ !!{proof} دے:
\[
\AxiomC{}
\RightLabel{$\pi_1$}
\Deduce$!B, !C, \Gamma \fCenter \Delta, !B$
\RightLabel{\LeftR{\land}}
\UnaryInf$!B \land !C, \Gamma \fCenter \Delta, !B$
\AxiomC{}
\RightLabel{$\pi_2$}
\Deduce$!B, !C, \Gamma \fCenter \Delta, !C$
\RightLabel{\LeftR{\land}}
\UnaryInf$!B \land !C, \Gamma \fCenter \Delta, !C$
\RightLabel{\RightR{\land}}
\BinaryInf$!B \land !C, \Gamma \fCenter \Delta, !B \land !C$
\DisplayProof
\]

د $!A \ident \lforall[x][!B(x)]$ دپاره، داسې !!a{constant} واخلئ چې په~$\lforall[x][!B(x)]$، $\Gamma$ يا~$\Delta$ کښې نۀ وي۔ د استقرا فرضيه د $!D = !B(c)$، $\Pi = \lforall[x][!B(x)], \Gamma$ او $\Lambda = \Delta$ دپاره تطبيق کړئ۔ دا د $!B(c),
\lforall[x][!B(x)], \Gamma \Sequent \Delta, !B(c)$ !!a{proof}~$\pi_1$ راکوي چې له هغې څخه دا ترلاسه کوو:
\[
\AxiomC{}
\RightLabel{$\pi_1$}
\Deduce$!B(c), \lforall[x][!B(x)], \Gamma \fCenter \Delta, !B(c)$
\RightLabel{\LeftR{\lforall}}
\UnaryInf$\lforall[x][!B(x)], \Gamma \fCenter \Delta, !B(c)$
\RightLabel{\RightR{\lforall}}
\UnaryInf$\lforall[x][!B(x)], \Gamma \fCenter \Delta, \lforall[x][!B(x)]$
\DisplayProof
\]
په $\RightR{\lforall}$ کښې د ځانګړي متغير شرط د~$c$ د ټاکلو د لارې له مخې پوره کېږي۔
\end{proof}

\begin{prob}
  د \olref{prop:G3c-axioms} ثبوت ته دوام ورکړئ: هغه صورتونه سر ته ورسوئ چې $!A \ident (!B \lif !C)$ او $!A \ident
  \lexists[x][B(x)]$ وي۔
\end{prob}

\end{document}
